\documentclass[10pt]{article}
\usepackage[margin=20mm]{geometry}
\usepackage{amsmath,amssymb,mathtools,fontspec,unicode-math,xurl,hyperref}
\IfFontExistsTF{DejaVu Serif}{\setmainfont{DejaVu Serif}}{\setmainfont{Latin Modern Roman}}
\setmathfont{latinmodern-math.otf}
\setlength{\parindent}{0pt}\setlength{\parskip}{6pt}\sloppy
\title{A Scope Correction in the Classification
of Non-uniform Graph Templates}
\author{AI-assisted technical review}
\date{}
\begin{document}
\maketitle
Abstract: This report examines the exact 2025 published version of Balister, Bollobás and Janson’s Consistent random vertex-orderings of graphs. The second sentence of Theorem 4.14 overstates the classification of an arbitrary, unreduced finite template. A three-vertex example disproves that intermediate assertion. The reduced-template classification and the infinite-blow-up consequence used in Theorem 1.3 are unaffected. This is a local scope defect, assessed as Minor. The moment, exchangeability and Bernoulli-polynomial arguments examined below yielded no further substantive issue. Scientific adjudication remains pending.\par
\section*{Paper identity and inspected scope}
The primary source is Paul N. Balister, Béla Bollobás and Svante Janson, Consistent random vertex-orderings of graphs, J. Eur. Math. Soc. 27 (2025), 2623–2652, doi: 10.4171/JEMS/1624 (https://doi.org/10.4171/JEMS/1624) [published]. The acquired publisher Version of Record is a 30-page PDF; its first page states a revision date of 7 November 2024. The publisher records publication on 27 March 2025. Here PDF page \(p\) corresponds to printed page \(2622+p\). The separately preserved arXiv:1506.03343v1, submitted 10 June 2015, is used only for the specific version comparison in Section 5 (\#sec:version).\par
All 30 pages of the published paper were read as extracted text. The claim-focused checks cover consistency and induced subgraphs; the constructions in Section 2; the rank representation in Section 3; the moment and template arguments leading to Theorems 1.3 and 1.5; the gluing and forest arguments in Section 5; and the Bernoulli calculations in Appendix A. Actual source-page checks are separately recorded with image hashes and performing roles. Standard external results such as de Finetti’s theorem, Ramsey’s theorem and the cited graph-limit equivalence were not independently reproved. No claim of a formal or exhaustive correctness certificate is made.\par
\section*{Definitions and the precise scope distinction}
A template \(T\) is a finite graph whose vertices are labelled full or empty. In \(T_\infty\), a vertex \(v\) is replaced by a countable blob \(W_v\), a clique if \(v\) is full and an independent set if \(v\) is empty. Edges between two blobs are complete or absent according to the corresponding edge of \(T\) [Retained latex: \textbackslash{}cite[Section 3, PDF p.\textasciitilde{}9]\{published\}].\par
A random ordering is consistent when the induced order distributions are preserved by every isomorphism between finite induced subgraphs. A template is uniform when its blow-up admits only the uniform consistent ordering. A template is reduced when it has neither adjacent full twins nor nonadjacent empty twins. Merging those twins preserves the isomorphism type of the countable blow-up, but need not preserve the finite underlying graph [Retained latex: \textbackslash{}cite[PDF p.\textasciitilde{}19]\{published\}].\par
\section*{A local scope defect in Theorem 4.14}
Problem 1 (The complete-or-empty assertion requires reduction).  The second sentence of Theorem 4.14, printed p. 2641 / PDF p. 19, asserts, for any non-uniform template \(G\), that \(G\) itself is complete or empty before deriving the classification of \(G_\infty\). The first sentence explicitly assumes reduction; the second does not. The unrestricted finite-template assertion is false.\par
\subsection*{A counterexample satisfying the literal hypotheses}
Let \[V(T)=\{a,b,c\},\qquad E(T)=\{\{a,b\}\},
 \qquad a,b,c\text{ all full}.\] This is a permitted finite template. It has one edge out of three possible edges, so its underlying graph is neither complete nor empty. By definition, \[T_\infty=C\sqcup D,\qquad C=W_a\cup W_b,\quad D=W_c,\] where \(C\) and \(D\) are countably infinite cliques.\par
Equation reference: eq:block. Use a fair coin to put all of \(C\) before all of \(D\), or vice versa, and use independent uniform random orders inside the two cliques. This is a measurable random total order. Its consistency can be verified for every finite induced subgraph, not merely for the three vertices used below. If that subgraph meets both cliques in \(r,s>0\) vertices, each of its allowed block orders has probability \[\frac{1}{2r!s!}.
 \] This depends only on the induced graph. It is invariant under internal permutations and, when applicable, interchange of the two components. If restriction removes a component, the remaining clique retains its uniform internal order. Thus every induced-subgraph isomorphism preserves the law.\par
Choose distinct \(a_1,a_2\in C\) and \(c_1\in D\). The block ordering gives \[\Pr(a_1<c_1<a_2)=0,\] whereas the uniform order on these three labelled vertices gives \(1/6\). Therefore \(T\) is non-uniform and satisfies every literal assumption of the second sentence. The claimed finite-graph conclusion fails. An exact enumeration of the six orders, retained in the evidence files, gives four probabilities \(1/4\) and two probabilities \(0\), summing to one. The enumeration checks this finite calculation; the preceding argument proves consistency for all finite induced subgraphs.\par
\subsection*{Counterevidence and supported correction}
The preceding paragraph of the paper says that it suffices to consider reduced templates. This makes the intended reduced-representative interpretation clear. It does not make the independently re-quantified unrestricted assertion true. In the example, \(a,b\) are adjacent full twins: the counterexample is deliberately not reduced and does not challenge the theorem’s first sentence.\par
Merging \(a,b\) produces a reduced template \(R\) consisting of two full vertices with no edge. Then \(R_\infty\cong T_\infty\), exactly as the reduction argument requires. A supported replacement for the second sentence is:\par
> For any non-uniform template \(T\), its blow-up \(T_\infty\) is a disjoint union of cliques or a complete multipartite graph.\par
Indeed, reduce \(T\) to \(R\) without changing the blow-up, apply the first sentence and its full/empty label restrictions to \(R\), and pass the resulting classification through \(R_\infty\cong T_\infty\). Alternatively, explicitly attribute the complete-or-empty assertion to the reduced representative.\par
\subsection*{Dependency and severity}
The proof of Theorem 1.3 on PDF p. 20 uses only the valid blow-up conclusion. Its witness graph is neither a disjoint union of cliques nor a complete multipartite graph and is an induced subgraph of the constructed blow-up. Both classes are hereditary, so the valid blow-up classification already gives the necessary contradiction. The subsequent use of Lemma 4.15 also does not require the false assertion about the unreduced finite template.\par
The assessment is Minor: a mathematical statement has an explicit counterexample, but the inspected central implication and classification of infinite blow-ups survive. The grading rests on that dependency analysis, not on the shortness of the correction. The counterexample is not evidence against Theorem 1.3, the reduced classification, or the paper’s general programme.\par
\section*{Other inspected reasoning}
\subsection*{Moment and exchangeability arguments}
The rank variables in Section 3 require more than marginal uniformity. Finite permutations within blobs give independent uniform internal orders; the tail-invariance argument then makes their rank labels independent of the random distribution functions. This supports the random-argument evaluations used in Section 4. Atoms are explicitly accommodated by comparing the relevant ranks; they are not silently discarded.\par
In Lemma 4.10, with the right-continuous inverse interpreted as in (3.8), \[\mathbb E X^2+\mathbb E Y^2
 =\frac23+\int_0^1(F(x)-x)^2\,dx.\] The equality case, monotonicity and right continuity identify the uniform distribution. This identity supports the equality arguments in Lemmas 4.11 and 4.15. The conditional expectations and Cauchy–Schwarz step in Lemma 4.12 were checked against the source equations (4.2)–(4.11); no additional substantive defect was established.\par
For Theorem 5.1, two independent uniform blocks supplied by the induction hypothesis force the de Finetti mixing measure to be deterministic: \(\mathbb E\mu(S)^2=(\mathbb E\mu(S))^2\) on a countable determining family. Conditioning on a positive-probability event in the fixed base ordering preserves exchangeability; the already deterministic empirical limit again forces the same mixing law. This justifies independence from the base ordering. The constructions for Theorem 1.6 require only vertex gluing and heredity. In Theorem 1.8, the finite pigeonhole step is supplemented by heredity to descend from an unbounded sequence of leaf counts to every required finite count. The edgeless and single-edge boundary cases were checked.\par
\subsection*{Non-uniform constructions and low-dimensional checks}
For Lemma 2.3, fixing an edge and imposing the signed sum modulo one gives Haar measure on an affine subtorus. Deleting any coordinate yields product uniform measure because every coefficient is \(1\) or \(-1\). Averaging over edges makes the law invariant under graph automorphisms. Proper subsets remain uniformly ordered.\par
The two Fourier calculations in Lemma A.2 reproduce its Bernoulli-polynomial coefficients. In the nonregular case, \[|E(G)|(P_{1,2}-P_{2,1})
 =(d_1-d_2)\frac{(-1)^n}{(n-1)!}\binom n2 B_{n-1}(\alpha).\] For \(n=3,\alpha=0\), the edge-conditioned difference is \(-1/4\), also obtained by an elementary area integral. In the regular case \(n=4\), the \(B_3\) coefficient cancels and \[|E(G)|(P_{1,2,3}-P_{2,1,3})=-B_2(\alpha).\] Thus \(\alpha=0\) gives \(-1/24\) for \(C_4\) and \(-1/12\) for \(2K_2\) after dividing by the respective numbers of edges. These are analytic checks, not simulation. The general argument uses linear independence of different-degree Bernoulli polynomials. Appendix A’s actual \(\omega=-2\pi it\) convention and inclusive \(n\ge3,n\ge4\) bounds were checked visually; extraction artefacts do not constitute findings.\par
\section*{Version comparison and remaining limits}
The preserved arXiv v1, Theorem 4.14 on its PDF p. 19, gives the reduced classification followed directly by the weaker, correct blow-up consequence [arxiv]. The extra unrestricted finite-template clause occurs in the reviewed publisher PDF. This identifies a difference between two acquired versions; it does not identify who introduced it or whether intermediate versions contain it.\par
A bounded title-plus-erratum search and inspection of the publisher and arXiv records did not identify a separate correction. This is not an assertion that none exists. The acquisition manifest also records bounded searches for supplements and paper-specific code, with none located. Their absence is not a defect in this theoretical paper. No author code was executed and no empirical reproduction is claimed. The cited external theory and every historical version have not been exhaustively checked.\par
\section*{Conclusion}
Dimension | Recorded state\par
Problem 1 | Minor: an unrestricted finite-template assertion fails.\par
Repair | The weaker blow-up consequence is supported by reduction; the inspected use in Theorem 1.3 is preserved.\par
Other checked claims | No further substantive issue established within the stated scope.\par
Scientific decision | Pending; no independent acceptance is claimed.\par
Author adoption | Unknown; no author contact or adoption request was made.\par
Remaining work | Independent focused acceptance of this report and evidence.\par
\section*{Supporting checks}
The main conversation performed the scientific assessment and prepared this report. Their actual source-page viewing is attributed individually. These are AI checks, not human reviews or the later independent acceptance stage.\par
Human review records are empty.\par
9 P. N. Balister, B. Bollobás and S. Janson, Consistent random vertex-orderings of graphs, J. Eur. Math. Soc. 27 (2025), 2623–2652. doi:10.4171/JEMS/1624 (https://doi.org/10.4171/JEMS/1624). Exact acquired publisher Version of Record. P. Balister, B. Bollobás and S. Janson, Consistent random vertex-orderings of graphs, arXiv:1506.03343v1 (https://arxiv.org/abs/1506.03343v1), 10 June 2015. Separately acquired author manuscript.\par
\end{document}
